% SGA 3, Expose I, Loop 1 English reconstruction.
% Scope: Section 4 opening through Definition 4.3.3.1, stopping immediately
% before Section 4.4 (Diagonalizable groups).
% Authority: Polo--Gille Exp1-13oct24.pdf, component-local pages 21--23
% (combined-reader pages 31--33). Force-OCR and jcreinhold are drafting and
% comparison witnesses only.

% Source page: Exp1 component-local p. 21; combined-reader p. 31.
% Re-edition printed p. 21; original-edition marginal pp. 23--24 also occur
% on this physical page.
\section*{\SGAInlineRefTarget{sga3:i:section:4}4. 概形范畴中的代数结构}
\addcontentsline{toc}{section}{4. 概形范畴中的代数结构}

在没有歧义时, 我们允许采用下述术语滥用. 用 \(T\) 表示
\(\mathscr C_{/S}\) 的对象 \(T\xrightarrow{f}S\), 而将 \(f\)
(即 \(T\) 的结构态射) 视为不言明的给定数据; 并将
\(\mathscr C\) 等同于 \(\widehat{\mathscr C}\) 的一个子范畴.
根据前几节所述的相容性, 这些等同不会造成问题.

我们也按照下述方式简化术语: 将 \((\mathbf{Sch})\)-群也称为
\emph{群概形}; 将 \((\mathbf{Sch})_{/S}\)-群称为
\emph{\(S\) 上的群概形}, \emph{\(S\)-群概形} 或 \emph{\(S\)-群};
当 \(S\) 是环 \(A\) 的谱时, 也称之为 \emph{\(A\)-群}.

\subsection*{\SGAInlineRefTarget{sga3:i:subsection:4-1}4.1. 常概形}

概形范畴满足 \SGARefLink{sga3:i:subsection:1:8}{1.8} 中要求的条件.
因此可以在其中定义常对象. 对每个集合 \(E\), 有一个概形
\(E_{\mathbb Z}\); 对每个概形 \(S\), 有一个 \(S\)-概形
\[
  E_S=(E_{\mathbb Z})_S.
\]
回顾一下 (参见\emph{上述文献}), 对每个 \(S\)-概形 \(T\),
\[
  \operatorname{Hom}_S(T,E_S)
\]
是从 \(T\) 的底拓扑空间到 \(E\) 的局部常值映射组成的集合.

函子 \(E\mapsto E_S\) 与有限逆极限可交换. 特别地, 若 \(G\) 是群,
则 \(G_S\) 是 \(S\)-群概形; 若 \(O\) 是环, 则 \(O_S\) 是
\(S\)-环概形; 其他代数结构亦然.

\subsection*{\SGAInlineRefTarget{sga3:i:subsection:4-2}4.2. 仿射 \(S\)-群}

先回顾关于仿射 \(S\)-概形的一些事实 (EGA~II, 第~1~节).
若 \(S\) 的每个仿射开子集的逆像都是仿射的, 就称 \(S\)-概形
\(T\) \emph{在 \(S\) 上仿射}. 此时 \(\mathcal O_S\)-代数
\(f_*\mathcal O_T\) 是拟凝聚的, 记为 \(\mathscr A(T)\), 其中 \(f\)
表示 \(T\) 的结构态射. 反之, 对每个拟凝聚 \(\mathcal O_S\)-代数
\(\mathscr A\), 都可以关联一个在 \(S\) 上仿射的 \(S\)-概形,
记为 \(\operatorname{Spec}(\mathscr A)\). 函子
\[
  T\longmapsto\mathscr A(T),
  \qquad
  \mathscr A\longmapsto\operatorname{Spec}(\mathscr A)
\]
是在 \(S\) 上仿射的 \(S\)-概形范畴与拟凝聚 \(\mathcal O_S\)-代数范畴
的反范畴之间的互为拟逆的等价.

由此可知, 在 \(S\) 上仿射的 \(S\)-概形 \(T\) 上给定代数结构,
等价于在拟凝聚 \(\mathcal O_S\)-代数范畴的反范畴中, 在
\(\mathscr A(T)\) 上给定相应的结构. 特别地, 若 \(G\) 是在 \(S\)
上仿射的 \(S\)-群, 则 \(\mathscr A(G)\) 带有增广
\(\mathcal O_S\)-双代数结构: 有两个 \(\mathcal O_S\)-代数态射
\[
  \Delta:\mathscr A(G)\longrightarrow
    \mathscr A(G)\otimes_{\mathcal O_S}\mathscr A(G),
  \qquad
  \varepsilon:\mathscr A(G)\longrightarrow\mathcal O_S,
\]
它们分别对应于 \(S\)-概形态射
\[
  \pi:G\times G\longrightarrow G,
  \qquad
  e_G:S\longrightarrow G.
\]
映射 \(\Delta\) 与 \(\varepsilon\) 满足以下条件, 这些条件表达了
\(G\) 是 \(S\)-幺半群.
\footnote{编者注: 当然, 取逆态射 \(G\to G\) 诱导一个
\(\mathcal O_S\)-代数态射 \(\tau:\mathscr A(G)\to\mathscr A(G)\),
它与 \(\Delta\), \(\varepsilon\) 一起使 \(\mathscr A(G)\) 成为
Hopf \(\mathcal O_S\)-代数.}

% Source page: Exp1 component-local p. 22; combined-reader p. 32.
% Re-edition printed p. 22; original-edition marginal p. 25 occurs below.
\medskip
\noindent\textbf{(HA 1)}
\(\Delta\) 是余结合的: 下图交换.

\begin{figure}[H]
  \centering
  \begin{tikzcd}[column sep=5.5em,row sep=3.8em]
  &
  \mathscr A(G)\otimes_{\mathcal O_S}\mathscr A(G)
    \arrow[dr,"\operatorname{Id}\otimes\Delta"]
  &
  \\
  \mathscr A(G)
    \arrow[ur,"\Delta"]
    \arrow[dr,"\Delta"']
  &&
  \mathscr A(G)\otimes_{\mathcal O_S}\mathscr A(G)
    \otimes_{\mathcal O_S}\mathscr A(G)
  \\
  &
  \mathscr A(G)\otimes_{\mathcal O_S}\mathscr A(G)
    \arrow[ur,"\Delta\otimes\operatorname{Id}"']
  &
  {} .
\end{tikzcd}

  \phantomsection\label{fig:sga3-exp1-p022-hopf-coassociativity}
\end{figure}

\noindent\textbf{(HA 2)}
\(\Delta\) 与 \(\varepsilon\) 相容: 以下两个复合均为恒等映射:
\[
  \mathscr A(G)
    \xrightarrow{\Delta}
  \mathscr A(G)\otimes_{\mathcal O_S}\mathscr A(G)
    \xrightarrow{\operatorname{Id}\otimes\varepsilon}
  \mathscr A(G)\otimes_{\mathcal O_S}\mathcal O_S
    \xrightarrow{\ \sim\ }
  \mathscr A(G),
\]
\[
  \mathscr A(G)
    \xrightarrow{\Delta}
  \mathscr A(G)\otimes_{\mathcal O_S}\mathscr A(G)
    \xrightarrow{\varepsilon\otimes\operatorname{Id}}
  \mathcal O_S\otimes_{\mathcal O_S}\mathscr A(G)
    \xrightarrow{\ \sim\ }
  \mathscr A(G).
\]

借此机会再强调一次, 根据 \(S\)-群结构的定义, 要在 \(S\) 上仿射的
\(S\)-概形 \(G\) 上给定这样的结构, 无须对 \(\mathscr A(G)\) 作任何验证;
只须对每个 \(S\) 上的 \(S'\), 在 \(G(S')\) 上赋予关于 \(S'\)
具有函子性的群结构. 作相应变通后, 这个观察也适用于态射.

\subsection*{\SGAInlineRefTarget{sga3:i:subsection:4-3}4.3. 群 \(\mathbb G_a\) 与 \(\mathbb G_m\). 环
\(\mathbf O\)}

\subsubsection*{4.3.1.}

令 \(\mathbb G_a\) 为从 \((\mathbf{Sch})^{\circ}\) 到
\(\mathbf{Set}\) 的函子, 定义为
\[
  \mathbb G_a(S)=\Gamma(S,\mathcal O_S),
\]
并赋予由环 \(\Gamma(S,\mathcal O_S)\) 的加法群结构定义的
\(\widehat{(\mathbf{Sch})}\)-群结构. 它可由一个仿射概形表示,
这个概形也记为 \(\mathbb G_a\), 因而是一个群概形:
\[
  \mathbb G_a=\operatorname{Spec}\mathbb Z[T].
\]
事实上,
\[
  \mathbb G_a(S)
    =\operatorname{Hom}(S,\mathbb G_a)
    =\operatorname{Hom}_{\mathrm{Alg.}}
       \bigl(\mathbb Z[T],\Gamma(S,\mathcal O_S)\bigr)
    \simeq\Gamma(S,\mathcal O_S).
\]

因此, 对每个概形 \(S\), 都有一个在 \(S\) 上仿射的 \(S\)-群,
记为 \(\mathbb G_{a,S}\), 它对应于 \(\mathcal O_S\)-双代数
\(\mathcal O_S[T]\), 其对角映射和增广由下式定义:
\[
  \Delta(T)=T\otimes1+1\otimes T,
  \qquad
  \varepsilon(T)=0.
\]

\subsubsection*{4.3.2.}

令 \(\mathbb G_m\) 为从 \((\mathbf{Sch})^{\circ}\) 到
\(\mathbf{Set}\) 的函子, 定义为
\[
  \mathbb G_m(S)=\Gamma(S,\mathcal O_S)^{\times},
\]
% Source page: Exp1 component-local p. 23; combined-reader p. 33.
% Re-edition printed p. 23; original-edition marginal p. 26 occurs below.
其中 \(\Gamma(S,\mathcal O_S)^{\times}\) 表示环 \(\Gamma(S,\mathcal O_S)\)
的可逆元组成的乘法群, 并赋予其自然的 \(\widehat{(\mathbf{Sch})}\)-群结构.
它可由一个仿射概形表示, 记为 \(\mathbb G_m\):
\[
  \mathbb G_m
    =\operatorname{Spec}\mathbb Z[T,T^{-1}]
    =\operatorname{Spec}\mathbb Z[\mathbb Z],
\]
其中 \(\mathbb Z[\mathbb Z]\) 表示 \(\mathbb Z\) 在环 \(\mathbb Z\)
上的群代数. 事实上,
\[
  \mathbb G_m(S)
    =\operatorname{Hom}_{\mathrm{Alg.}}
       \bigl(\mathbb Z[T,T^{-1}],\Gamma(S,\mathcal O_S)\bigr)
    \simeq\Gamma(S,\mathcal O_S)^{\times}.
\]

因此, 对每个概形 \(S\), 都有一个在 \(S\) 上仿射的 \(S\)-群,
记为 \(\mathbb G_{m,S}\), 它对应于 \(\mathcal O_S\)-代数
\(\mathcal O_S[\mathbb Z]\), 其对角映射和增广由下式定义:
\[
  \Delta(x)=x\otimes x,
  \qquad
  \varepsilon(x)=1,
  \qquad x\in\mathbb Z.
\]

\subsubsection*{4.3.3.}

令 \(\mathbf O\) 为赋予了 \(\widehat{(\mathbf{Sch})}\)-环结构的函子
\(\mathbb G_a\). 它由概形 \(\operatorname{Spec}\mathbb Z[T]\) 表示;
在将该概形视为带有环概形结构时, 我们将其记为 \(\mathbf O\).

因此, 对每个概形 \(S\),
\[
  \mathbf O_S
    =S\times_{\operatorname{Spec}\mathbb Z}
       \operatorname{Spec}\mathbb Z[T]
    =\operatorname{Spec}(\mathcal O_S[T])
\]
是在 \(S\) 上仿射的 \(S\)-环概形. (注意: 在 EGA~II,
1.7.13 中, \(\mathbf O_S\) 记为 \(S[T]\).)

\medskip
\noindent\SGAInlineRefTarget{sga3:i:paragraph:4-3-3-1}\textbf{4.3.3.1.}
\footnote{编者注: 本段为新增内容; 后文将会用到 (参见 II,
\SGARefLink{sga3:i:subsection:1-3}{1.3}).}
对 \(\widehat{(\mathbf{Sch})}\) 的每个对象 \(X\),
\[
  \mathbf O(X)=\operatorname{Hom}(X,\mathbf O)
\]
都带有关于 \(X\) 具有函子性的环结构. 特别地, 若 \(X'\) 是一个概形,
并给定态射 \(x:X'\to X\) 和 \(f:X\to\mathbf O\)
(即 \(f\in\mathbf O(X)\)), 则
\[
  f(x)=f\circ x
    \in\mathbf O(X')
    =\Gamma(X',\mathcal O_{X'}).
\]

\medskip
\noindent\textbf{定义.}
设 \(\pi:M\to X\) 是 \(\widehat{(\mathbf{Sch})}\) 中的态射, 并令
\[
  \mathbf O_X=\mathbf O\times X.
\]
若对每个 \(X\)-概形 \(X'\), 集合
\[
  \operatorname{Hom}_X(X',M)
\]
都赋予了关于 \(X'\) 具有函子性的 \(\mathbf O(X')\)-模结构,
则称 \(M\) 为 \(\mathbf O_X\)-模.

这等价于给定一个 \(X\)-交换群运算
\[
  \mu:M\times_X M\longrightarrow M
\]
和一个\emph{外运算}
\[
  \mathbf O\times M
    =\mathbf O_X\times_X M
    \longrightarrow M,
  \qquad
  (f,m)\longmapsto f\cdot m,
\]
后者是 \(X\)-态射, 即满足 \(\pi(f\cdot m)=\pi(m)\), 且对每个
\(x\in X(S)\), 它赋予
\[
  M(x)=\{m\in M(S)\mid\pi(m)=x\}
\]
以 \(\mathbf O(S)\)-模结构.

在此情形下, 对每个不一定可表的
\(Y\in\operatorname{Ob}\bigl(\widehat{(\mathbf{Sch})}_{/X}\bigr)\),
\[
  \operatorname{Hom}_X(Y,M)=\Gamma(M_Y/Y)
\]
都是一个关于 \(Y\) 具有函子性的 \(\mathbf O(Y)\)-模.

% End of assigned component. Section 4.4 (Diagonalizable groups) begins next,
% on Exp1 component-local p. 23 / combined-reader p. 33.
% Translation statistics: source_lines=288; output_lines=267; inline_math=148; display_math=24; align_environments=0.
% Review: complete assigned component; pending_TODO=0; standalone_compile=blocked_missing_ctex.sty; integrated_compile=not_run.
